\documentclass[10pt,letterpaper]{article}
\usepackage[utf8]{inputenc}
\usepackage[spanish]{babel}
\usepackage{enumitem}
\usepackage{hyperref}
\usepackage{xcolor}
\usepackage{fontawesome5}
\usepackage{parskip}
\usepackage[margin=0.75in]{geometry}
\usepackage{setspace}

\hypersetup{
    colorlinks=true,
    linkcolor=black,
    urlcolor=black,
    citecolor=black,
    filecolor=black
}

\setlist[itemize]{leftmargin=*,nosep,after=\vspace{0.1cm}}
\setlist[enumerate]{leftmargin=*,nosep}
\setlength{\parskip}{0.3cm}

\begin{document}

% Header
\begin{center}
    {\Huge \textbf{Bruno Camilo Henríquez Cano}} \\[0.2cm]
    {\Large Ingeniero en Maquinaria Industrial $\cdot$ Especialista en Sistemas Mecánicos}
\end{center}

\vspace{0.3cm}

% Contact
\noindent
\begin{minipage}[t]{0.5\textwidth}
    \faMapMarker\ Santiago, Chile \\
    \faPhone\ +56 9 2815 6161 \\
    \faEnvelope\ \href{mailto:bruno.henriquez.1993@gmail.com}{bruno.henriquez.1993@gmail.com}
\end{minipage}
\begin{minipage}[t]{0.5\textwidth}
    \faLinkedin\ \href{https://linkedin.com/in/bruno-henriquezcano}{linkedin.com/in/bruno-henriquezcano} \\
    \faGithub\ \href{https://github.com/brunochc}{github.com/brunochc} \\
    \faGlobe\ \href{https://bchc.tech}{bchc.tech} \\
    \faCar\ Licencia de Conducir Clase B
\end{minipage}

\vspace{0.4cm}

% Professional Profile
\noindent{\Large \textbf{PERFIL PROFESIONAL}}

\vspace{0.2cm}

\noindent Ingeniero especializado en diagnóstico, reparación y mantenimiento de maquinaria industrial y equipos pesados. Más de 7 años de experiencia práctica resolviendo fallas complejas en sistemas hidráulicos, mecánicos y electrónicos. Capacidad demostrada para interpretar planos técnicos, realizar diagnósticos precisos y ejecutar reparaciones de alta calidad. Experiencia con flotas de +100 equipos mineros y +400 vehículos de transporte. Conocimiento sólido de SAP-PM, mantenimiento preventivo y correctivo, y gestión de repuestos críticos.

\vspace{0.4cm}

% Work Experience
\noindent{\Large \textbf{EXPERIENCIA LABORAL}}

\vspace{0.3cm}

\noindent \textbf{Secretario de Planificación Técnica} \\
\textit{Consorcio Acciona Ossa Pizzarotti (Codelco) — Calama — Diciembre 2023 – Noviembre 2024}
\begin{itemize}
    \item Gestioné el mantenimiento técnico de una flota de más de 100 equipos mineros pesados (palas hidráulicas, camiones de extracción, perforadoras), asegurando disponibilidad operativa continua.
    \item Realicé diagnósticos de fallas complejas en sistemas hidráulicos, eléctricos y mecánicos, identificando causas raíz y coordinando reparaciones efectivas.
    \item Implementé planes de mantenimiento preventivo basados en horas de operación y condiciones de trabajo, reduciendo fallas inesperadas.
    \item Coordiné con proveedores la adquisición de repuestos críticos (componentes hidráulicos, sistemas de transmisión, filtros especializados), asegurando disponibilidad y calidad.
    \item Utilicé SAP-PM para registro de órdenes de trabajo, historial de mantenimiento y trazabilidad de componentes.
\end{itemize}

\vspace{0.2cm}

\noindent \textbf{Coordinador de Sistemas y Análisis de Flota} \\
\textit{STP Santiago S.A — Maipú — Junio 2021 – Agosto 2023}
\begin{itemize}
    \item Supervisé el estado técnico de una flota de más de 400 buses de transporte público, gestionando programas de mantenimiento preventivo y correctivo.
    \item Diagnostiqué fallas recurrentes en sistemas de transmisión, frenos, suspensión y motor, implementando soluciones que mejoraron la disponibilidad en 15\%.
    \item Analicé datos de telemetría (OBD-II) para identificar patrones de desgaste y predecir fallas antes de que ocurrieran.
    \item Coordiné con el taller mecánico la priorización de reparaciones según criticidad operativa.
    \item Generé reportes técnicos de confiabilidad (MTBF, MTTR) para la gerencia de mantenimiento.
\end{itemize}

\vspace{0.2cm}

\noindent \textbf{Técnico Automotriz Independiente} \\
\textit{Autónomo — Paine, Chile — 2014 – 2020}
\begin{itemize}
    \item Diagnostiqué y reparé fallas complejas en más de 30 vehículos mensuales, especializándome en casos "sin solución" de otros talleres.
    \item Ejecuté reparaciones en sistemas de inyección electrónica, transmisión automática, diagnóstico de redes CAN, suspensión y dirección.
    \item Utilicé herramientas de diagnóstico avanzado (scanner OBD-II, multímetro automotriz, osciloscopio) para identificar fallas eléctricas y electrónicas.
    \item Mantuve un 95\% de satisfacción del cliente mediante diagnósticos precisos y comunicación transparente sobre alcance de reparaciones.
    \item Gestioné relaciones con proveedores de repuestos, asegurando calidad y tiempos de entrega competitivos.
\end{itemize}

\vspace{0.4cm}

% Technical Skills
\noindent{\Large \textbf{HABILIDADES TÉCNICAS}}
\vspace{0.3cm}

\noindent
\begin{minipage}[t]{0.48\textwidth}
\textbf{Diagnóstico y Reparación:}
\begin{itemize}
    \item Sistemas Hidráulicos (circuitos, válvulas, bombas)
    \item Sistemas Mecánicos (transmisión, motor, suspensión)
    \item Sistemas Eléctricos y Electrónicos
    \item Lectura de planos técnicos y esquemáticos
    \item Diagnóstico con OBD-II, multímetro, osciloscopio
\end{itemize}

\vspace{0.2cm}

\textbf{Mantenimiento:}
\begin{itemize}
    \item Mantenimiento Preventivo y Predictivo
    \item Análisis de Causa Raíz (RCA)
    \item Gestión de Repuestos Críticos
    \item SAP-PM (gestión de órdenes de trabajo)
    \item KPIs de Mantenimiento (MTBF, MTTR, OEE)
\end{itemize}
\end{minipage}
\hfill
\begin{minipage}[t]{0.48\textwidth}
\textbf{Maquinaria y Equipos:}
\begin{itemize}
    \item Equipos Mineros (palas, camiones, perforadoras)
    \item Vehículos de Transporte (buses, camiones)
    \item Maquinaria Agrícola y de Construcción
    \item Sistemas de Alta Presión Hidráulica
    \item Motores Diesel y Gasolina
\end{itemize}

\vspace{0.2cm}

\textbf{Herramientas de Gestión:}
\begin{itemize}
    \item SAP-PM (Planificación de Mantenimiento)
    \item Microsoft Excel (análisis de datos técnicos)
    \item Software de diagnóstico automotriz
    \item Sistemas de telemetría vehicular
\end{itemize}
\end{minipage}

\vspace{0.3cm}

% Education
\vspace{0.4cm}
\noindent{\Large \textbf{EDUCACIÓN}}

\vspace{0.3cm}

\noindent \textbf{Ingeniería en Maquinaria, Sistemas Automotrices y Electrónicos} \\
\textit{INACAP — Marzo 2014 – Octubre 2019}

\vspace{0.2cm}

\noindent \textbf{Ingeniería Civil en Computación (2 años completados)} \\
\textit{Universidad Tecnológica Metropolitana — Marzo 2012 – Marzo 2014}

\vspace{0.4cm}

% Additional Training
\noindent{\Large \textbf{FORMACIÓN COMPLEMENTARIA}}

\vspace{0.3cm}

\begin{itemize}
    \item Diagnóstico Electrónico Automotriz Avanzado
    \item Sistemas Hidráulicos Industriales
    \item Gestión de Mantenimiento con SAP-PM
    \item Análisis de Fallas en Maquinaria Pesada
\end{itemize}

\vspace{0.4cm}

% Languages
\noindent{\Large \textbf{IDIOMAS}}

\vspace{0.3cm}

\noindent \textbf{Español:} Nativo \\
\noindent \textbf{Inglés:} Intermedio (lectura de manuales técnicos y documentación)

\end{document}
